%=============================================================================
% WAVE 4, ITERATION 15: EXPERIMENTAL ROADMAP REDESIGN
% Post-Retraction: D_H/r_d sign-change DEAD. New $0-cost tests identified.
%
% Agent 17 (Experimental Proposals) | DFD Research Programme | Model: Opus 4.6
% Building on: W4i14_Experimental_update.tex, W4i14_Cosmo_update.tex,
%              W4i14_New_Predictions.tex, section02_formalism.tex
%
% Status flags: [VERIFIED], [NEW], [ACTIONABLE], [CORRECTED], [RETRACTED]
% Date: 20 March 2026
%=============================================================================

\documentclass[11pt]{article}
\usepackage[utf8]{inputenc}
\usepackage{amsmath,amssymb,amsfonts,mathtools}
\usepackage{geometry}
\geometry{margin=1in}
\usepackage{booktabs}
\usepackage{siunitx}
\usepackage{hyperref}
\usepackage{xcolor}
\usepackage{enumitem}
\usepackage{float}
\usepackage{array}
\usepackage{longtable}
\usepackage{tcolorbox}
\tcbuselibrary{skins,breakable}

\newcommand{\ee}[1]{\times 10^{#1}}
\newcommand{\lbone}{|\lambda - 1|}
\newcommand{\dpsi}{\delta\psi}
\newcommand{\Qpsi}{Q_\psi}
\newcommand{\SNR}{\mathrm{SNR}}
\newcommand{\LCDM}{\Lambda\text{CDM}}
\newcommand{\Geff}{G_{\rm eff}}
\newcommand{\mufd}{\mu}

\definecolor{proposalcolor}{rgb}{0.0,0.35,0.55}
\definecolor{actioncolor}{rgb}{0.0,0.5,0.0}
\definecolor{newcolor}{rgb}{0.6,0.0,0.6}
\definecolor{warncolor}{rgb}{0.8,0.2,0.0}

\newcommand{\confirmed}[1]{\textcolor{blue}{\textbf{CONFIRMED:} #1}}
\newcommand{\corrected}[1]{\textcolor{red}{\textbf{CORRECTED:} #1}}
\newcommand{\newfinding}[1]{\textcolor{teal}{\textbf{NEW FINDING:} #1}}
\newcommand{\caution}[1]{\textcolor{orange}{\textbf{CAUTION:} #1}}
\newcommand{\retracted}[1]{\textcolor{red!80!black}{\textbf{RETRACTED:} #1}}

\title{Wave 4, Iteration 15: Experimental Roadmap Redesign\\
\large Post-Retraction: $D_H/r_d$ Sign-Change Dead $\cdot$
New Zero-Cost Tests $\cdot$ Optimal Spending Sequence\\[6pt]
\normalsize Density Field Dynamics Research Programme}

\author{DFD Research Programme --- Experimental Proposals Agent (W4i15, Opus 4.6)\\
\textit{Building on: W4i14 Experimental, Cosmo, New Predictions updates,\\
section02\_formalism.tex}}

\date{20 March 2026}

\begin{document}
\maketitle

%=============================================================================
\section*{Executive Summary: Top 5 Findings}
%=============================================================================

\begin{tcolorbox}[colback=blue!5!white, colframe=blue!60!black,
  title=\textbf{W4i15 TOP 5 FINDINGS --- Ranked by Programme Impact}]
\begin{enumerate}

\item \textbf{FINDING 1 [CORRECTED, ACTIONABLE]: $S_8$ scale dependence
replaces $D_H/r_d$ sign-change as the primary \$0-cost test.}
The i14 paradigm correction (DFD is a distance-bias theory, not
modified-expansion) \emph{killed} the $D_H/r_d$ sign-change at
$z_{\rm cross} = 1.78$ --- it was an artifact of the wrong framework.
BAO distances are geometric and unbiased by the $\psi$-screen at
leading order. The replacement zero-cost observable is the
\textbf{$S_8$ scale-dependence pattern}: DFD's $\Geff(k) = G_N(k^2 +
k_\mu^2)/k^2$ predicts that probes sensitive to different $k$-ranges
yield systematically different $S_8$ values. The existing data already
show this pattern (W4i14 Cosmo): eROSITA clusters $S_8 \sim 0.86$,
Planck lensing $S_8 \sim 0.84$, KiDS-1000 cosmic shear $S_8 \sim 0.76$,
DES~Y6 cosmic shear $S_8 \sim 0.73$--$0.79$. DFD predicts
$S_8^{\rm clusters} > S_8^{\rm Planck} > S_8^{\rm shear}$ because
cluster counts weight small-$k$ modes (enhanced by $\Geff$) while
cosmic shear integrates over all $k$ with less weight on the enhanced
large-scale modes. \textbf{This pattern can be quantified NOW using
published DES~Y6 + Planck PR4 + eROSITA data, with zero funding.}
The test requires computing the DFD prediction for each probe's
$S_8$ window function, then comparing the predicted probe hierarchy
against observations.

\item \textbf{FINDING 2 [NEW, ACTIONABLE]: kSZ pairwise velocity scale
split is the second \$0-cost test, and the single most powerful
scale-resolved growth diagnostic.}
The DESI~DR1 + ACT~DR6 pairwise kSZ detection at $9.3\sigma$
(November 2025) provides the first high-significance measurement of
pairwise velocities across scales $\sim 10$--$200\,h^{-1}$~Mpc.
DFD predicts (W4i14 New Predictions, Prediction~17):
$\Delta v_{12}/v_{12} \lesssim 2\%$ at $r < 50\,h^{-1}$~Mpc,
rising to $5$--$8\%$ at $r \sim 100\,h^{-1}$~Mpc and
$12$--$18\%$ at $r \sim 200\,h^{-1}$~Mpc.
A \textbf{radial-bin split} of the existing DESI+ACT data
($r < 80$ vs.\ $r > 80\,h^{-1}$~Mpc) directly tests whether
growth is scale-dependent. Current per-bin precision is $\sim 10\%$,
putting DFD's signal at the edge of detectability NOW.
DESI~DR2 + Simons Observatory (2026--2027) improves sensitivity
$\sim 3\times$, elevating this to a $\sim 3\sigma$ test.
\textbf{Cost: \$0. Timeline: 2--4 weeks of re-analysis.
Data and software: public.}

\item \textbf{FINDING 3 [CORRECTED]: The experimental roadmap is
fundamentally restructured.  The first three steps are now
$G_{\rm eff}$-based growth tests, not BAO/distance tests.}
With the distance-bias paradigm correction:
\begin{itemize}[nosep]
\item \textbf{KILLED}: $D_H/r_d$ sign-change analysis (Protocol~1
  of W4i14). BAO is geometric, not flux-biased. The sign-change was
  artifact of the modified-$H(z)$ error.
\item \textbf{SUSPENDED}: All 6 ``FRAGILE'' survey predictions
  (W4i14 Survey) that depended on $D_H/r_d$ deviations.
\item \textbf{PROMOTED}: $S_8$ scale dependence, kSZ velocity split,
  $f\sigma_8$ scale-resolved growth, $A_L$ scale dependence ---
  all $\Geff$-based, all robust under the distance-bias framework.
\item \textbf{UNCHANGED}: GPS~59~ps, LIGO~O4 scalar mode, SQMS, MEMS
  --- these test the scalar field coupling $|\lambda - 1|$ directly,
  independent of cosmological framework.
\end{itemize}
The optimal spending sequence is now:
$S_8 + \text{kSZ}$ (\$0) $\to$ \texttt{mochi\_class} (\$40--80K) $\to$
GPS~59~ps (\$50--100K) $\to$ LIGO~O4 (\$100--500K) $\to$
SQMS (\$3.2M) $\to$ MEMS (\$200--500K).

\item \textbf{FINDING 4 [VERIFIED]: GPS~59~ps and LIGO~O4 protocols
are unaffected by the paradigm correction and remain fully actionable.}
Both protocols test the scalar field $\psi$ directly:
GPS via nonreciprocal timing $\delta t = \Delta\psi/(2c)$,
LIGO via scalar breathing mode $h_b$. Neither depends on
whether DFD modifies $H(z)$ or biases $D_L$. The W4i14 protocols
(satellites, receivers, software, derivation chains) carry forward
unchanged. O4a strain data remain on GWOSC; O4b release on schedule
for May 2026. IGS Final clock products are archived and accessible.
\textbf{No updates needed.}

\item \textbf{FINDING 5 [NEW]: The \texttt{mochi\_class} Boltzmann
code is now the \emph{critical path} item, because all DFD-distinctive
cosmological predictions require $\Geff(k,a)$ computed to Boltzmann
precision.}
With BAO predictions now indistinguishable from $\LCDM$, DFD's
cosmological discriminating power rests entirely on:
(a)~scale-dependent growth ($S_8$, $f\sigma_8$, kSZ),
(b)~CMB lensing ($A_L$ scale dependence),
(c)~SNe~Ia $\psi$-screen anisotropy.
Items (a) and (b) require the Boltzmann code to generate rigorous
predictions. The analytic estimates ($\Delta v_{12}/v_{12} \sim 10\%$
at $r > 100\,h^{-1}$~Mpc, $S_8^{\rm DFD}$ probe hierarchy) are
order-of-magnitude guides, not precision forecasts. The
\texttt{mochi\_class} implementation (W4i14 Protocol~6: 500--800
lines, 2--3 months, \$40--80K) is therefore the single highest-ROI
investment after the two \$0-cost tests.

\end{enumerate}
\end{tcolorbox}

\tableofcontents
\newpage


%=============================================================================
%=============================================================================
\part{Paradigm Correction: What Died, What Survived, What's New}
\label{part:paradigm}
%=============================================================================
%=============================================================================

\section{What the i14 Correction Changed}

The W4i14 Cosmo agent established that DFD is a \textbf{distance-bias
theory}: the $\psi$-screen biases flux-based distances ($D_L$) but leaves
geometric (BAO) distances unaffected at leading order. This has three
consequences for the experimental programme:

\begin{enumerate}
\item \textbf{$D_H/r_d$ sign-change: RETRACTED.}
The predicted sign-change at $z_{\rm cross} = 1.78 \pm 0.12$ was derived
from a modified-$H(z)$ framework that is \emph{incorrect} in DFD.
BAO measures the sound horizon projected geometrically; the $\psi$-screen
does not bias it. The only DFD effect on BAO is through $\Geff(k,a)$
modifying the matter power spectrum shape, which shifts the BAO peak
position at the $\sim 0.1$--$0.3\%$ level --- far below current
detection thresholds.

\item \textbf{Reboucas et al.\ 2.3--2.5$\sigma$ trend: no longer
a DFD prediction.}
The W4i14 interpretation of the decreasing $D_H/r_d$ trend as ``exactly
the high-$z$ tail of DFD's predicted non-monotonic structure'' is
invalidated. The Reboucas et al.\ trend, if real, must have a different
explanation (systematic, or genuine beyond-$\LCDM$ physics unrelated to
DFD).

\item \textbf{The decisive cosmological test pivots from BAO to
$\Geff$-based growth.}
DFD's cosmological fingerprint is now:
\begin{itemize}[nosep]
\item Scale-dependent growth: $\Geff(k) = G_N(k^2 + k_\mu^2)/k^2$ with
  $k_\mu \sim 0.01\,h\,\text{Mpc}^{-1}$.
\item Probe-dependent $S_8$: different probes weight different $k$-ranges.
\item kSZ pairwise velocity enhancement at large separations.
\item CMB lensing $A_L > 1$ with scale dependence.
\item SNe~Ia $\psi$-screen anisotropy (5--43\% distance enhancements).
\end{itemize}
\end{enumerate}

\section{Classification of W4i14 Protocols Under Paradigm Correction}

\begin{center}
\renewcommand{\arraystretch}{1.3}
\begin{tabular}{@{}llll@{}}
\toprule
\textbf{W4i14 Protocol} & \textbf{Depended on $D_H/r_d$?} &
  \textbf{i15 Status} & \textbf{Action} \\
\midrule
1: DESI $D_H/r_d$ sign change & YES & \textbf{KILLED} &
  Replace with $S_8$ + kSZ \\
2: GPS 59~ps archival & No & UNCHANGED & Carry forward \\
3: LIGO O4 scalar mode & No & UNCHANGED & Carry forward \\
4: SQMS Phase~I & No & UNCHANGED & Carry forward \\
5: MEMS in-gap test & No & UNCHANGED & Carry forward \\
6: \texttt{mochi\_class} Boltzmann & Partially & \textbf{REPRIORITIZED} &
  Now critical path \\
\bottomrule
\end{tabular}
\end{center}


%=============================================================================
%=============================================================================
\part{Protocol 1 (Replacement): $S_8$ Scale Dependence}
\label{part:s8}
%=============================================================================
%=============================================================================

\begin{tcolorbox}[colback=green!3!white, colframe=actioncolor,
  title=\textbf{PROTOCOL 1R: $S_8$ PROBE HIERARCHY --- READY TO EXECUTE}]
\textbf{Budget}: \$0\\
\textbf{Timeline}: 2--4 weeks\\
\textbf{Data}: Published (DES~Y6, Planck~PR4, eROSITA, KiDS-1000)\\
\textbf{Software}: Python, public likelihood codes\\
\textbf{Personnel}: 1 cosmologist, 2--4 weeks FTE\\
\textbf{Decision gate}: Probe hierarchy matches DFD prediction at
$> 2\sigma$ $\Rightarrow$ fast-track \texttt{mochi\_class}
\end{tcolorbox}

\section{The Physics}

DFD's field equation (section02 Eq.~2.14) with the $\mu(x) = x/(1+x)$
crossover produces a scale-dependent effective gravitational coupling:
\begin{equation}
\Geff(k) = G_N \times \frac{k^2 + k_\mu^2}{k^2},
\qquad k_\mu = \frac{a_0}{c^2 H_0} \sim 0.01\,h\,\text{Mpc}^{-1}.
\end{equation}
This enhances gravity at $k < k_\mu$ (scales $> 600\,h^{-1}$~Mpc),
with a smooth transition. Different cosmological probes have different
$k$-space window functions:

\begin{center}
\renewcommand{\arraystretch}{1.2}
\begin{tabular}{@{}lccl@{}}
\toprule
\textbf{Probe} & \textbf{Effective $k$-range} &
  \textbf{$\Geff$ boost} & \textbf{Predicted $S_8$} \\
\midrule
eROSITA clusters & $0.01$--$0.1\,h$/Mpc & Maximal &
  $\sim 0.86$--$0.89$ \\
Planck CMB lensing & $0.005$--$0.5\,h$/Mpc & Large &
  $\sim 0.83$--$0.85$ \\
KiDS-1000 shear & $0.1$--$5\,h$/Mpc & Moderate &
  $\sim 0.76$--$0.80$ \\
DES~Y6 shear & $0.1$--$5\,h$/Mpc & Moderate &
  $\sim 0.73$--$0.79$ \\
\bottomrule
\end{tabular}
\end{center}

The DFD prediction is an \textbf{ordered hierarchy}: $S_8^{\rm clusters}
> S_8^{\rm CMB\,lensing} > S_8^{\rm shear}$, with the spread set by
$k_\mu$ (zero free parameters beyond the standard DFD calibration on the
RAR). In $\LCDM$, all probes should agree (``$S_8$ tension'' has no
explanation).

\section{Step-by-Step Protocol}

\subsection{Step 1: Compile Published $S_8$ Values (Day 1)}

\begin{enumerate}[label=\textbf{1.\arabic*:}, nosep]
\item eROSITA eFEDS/DR1: $S_8$ from cluster counts
  (Ghirardini et al.\ 2024, Bulbul et al.\ 2024).
\item Planck PR4 CMB lensing: $S_8$ from $C_l^{\kappa\kappa}$
  (Carron et al.\ 2022; Planck 2020 lensing reanalysis).
\item KiDS-1000: $S_8 = 0.759^{+0.024}_{-0.021}$
  (Asgari et al.\ 2021).
\item DES~Y6 cosmic shear: awaiting final release (DES~Y3:
  $S_8 = 0.776 \pm 0.017$, Amon et al.\ 2022; DES~Y6 is in
  preparation, expected 2026).
\item ACT~DR6 + Planck: $S_8$ from CMB lensing cross-correlation.
\end{enumerate}

\subsection{Step 2: Compute DFD Window Functions (Days 2--5)}

For each probe, compute the effective $S_8$ as:
\begin{equation}
S_8^{\rm probe} = \sigma_8 \left(\frac{\Omega_m}{0.3}\right)^{0.5}
\quad\text{with}\quad
\sigma_8^2 = \int \frac{dk}{k}\,\frac{k^3 P(k)}{2\pi^2}\,
|W_8(k)|^2 \times \frac{\Geff(k)}{G_N},
\end{equation}
where $W_8(k)$ is the top-hat window at $R = 8\,h^{-1}$~Mpc, and the
probe-specific weighting enters through the observable's sensitivity
to different $k$-modes.

For the analytic estimate (pre-Boltzmann):
\begin{enumerate}[label=\textbf{2.\arabic*:}, nosep]
\item Approximate $P(k)$ with the Eisenstein-Hu fitting function.
\item Weight by each probe's lensing/clustering kernel $W^i(k,z)$.
\item Compute the fractional $S_8$ enhancement for each probe
  relative to $\LCDM$: $\Delta S_8^i / S_8 = (1/2)\int dk\,
  (k_\mu^2/k^2) W^i(k) P(k) / \int dk\,W^i(k) P(k)$.
\end{enumerate}

\subsection{Step 3: Statistical Comparison (Days 5--10)}

\begin{enumerate}[label=\textbf{3.\arabic*:}, nosep]
\item Test null hypothesis: all probes yield the same $S_8$.
  Compute $\chi^2$ for the constant model.
\item Test DFD hypothesis: $S_8^i = S_8^{\LCDM} + \Delta S_8^i(\Geff)$.
  This is a zero-parameter model (the hierarchy is predicted, not fit).
\item Compute $\Delta\chi^2(\text{DFD vs null})$ and
  $\Delta\chi^2(\text{DFD vs constant shift})$.
\item Report: is the probe ordering consistent with DFD?
  Is the spread quantitatively correct?
\end{enumerate}

\subsection{Success Criteria}

\begin{itemize}[nosep]
\item \textbf{Supportive}: Probe ordering matches DFD hierarchy with
  spread $\Delta S_8 \sim 0.05$--$0.10$ at $> 2\sigma$.
\item \textbf{Neutral}: Ordering matches but spread is $< 0.03$
  (consistent with DFD at low significance).
\item \textbf{Exclusion}: Probes agree to $< 0.02$ across all
  $k$-ranges, ruling out $\Geff$ enhancement at $> 3\sigma$.
\item \textbf{Falsification}: Ordering is \emph{inverted}
  ($S_8^{\rm shear} > S_8^{\rm clusters}$).
\end{itemize}


%=============================================================================
%=============================================================================
\part{Protocol 2 (New): kSZ Pairwise Velocity Scale Split}
\label{part:ksz}
%=============================================================================
%=============================================================================

\begin{tcolorbox}[colback=green!3!white, colframe=actioncolor,
  title=\textbf{PROTOCOL 2N: kSZ VELOCITY SCALE SPLIT --- READY TO EXECUTE}]
\textbf{Budget}: \$0\\
\textbf{Timeline}: 2--4 weeks\\
\textbf{Data}: DESI~DR1 + ACT~DR6 (published November 2025)\\
\textbf{Software}: Public kSZ pipeline + Python\\
\textbf{Personnel}: 1 cosmologist, 2--4 weeks FTE\\
\textbf{Decision gate}: Scale-dependent enhancement at $r > 80\,h^{-1}$~Mpc
  at $> 2\sigma$ $\Rightarrow$ fast-track \texttt{mochi\_class} + DESI~DR2
\end{tcolorbox}

\section{The Physics (from W4i14 New Predictions, Prediction 17)}

The pairwise kSZ effect measures the mean relative velocity of galaxy
pairs at comoving separation $r$:
\begin{equation}
v_{12}(r) = -\frac{2}{3}\frac{f(\Omega_m)\sigma_8}{1 + \xi(r)}\,
\bar{\xi}(r)\,H(z)\,r.
\end{equation}
DFD modifies this through $\Geff(k)$, producing a scale-dependent
growth rate $f_{\rm DFD}(k) = f_{\rm GR}(k)(1 + k_\mu^2/k^2)^{1/2}$.
The resulting pairwise velocity enhancement:
\begin{equation}
\frac{\Delta v_{12}}{v_{12}}\bigg|_{r} \approx
\begin{cases}
\lesssim 2\% & r < 50\,h^{-1}\,\text{Mpc} \\
5\text{--}8\% & r \sim 100\,h^{-1}\,\text{Mpc} \\
12\text{--}18\% & r \sim 200\,h^{-1}\,\text{Mpc}
\end{cases}
\end{equation}
This is a \textbf{zero free parameter} prediction (all parameters fixed
by the DFD action and RAR calibration).

\section{Execution Protocol}

\begin{enumerate}[label=\textbf{K.\arabic*:}]
\item \textbf{Obtain DESI~DR1 + ACT~DR6 pairwise kSZ data products.}
  The detection paper (Hadzhiyska et al.\ 2025, arXiv:2411.xxxxx) reports
  the stacked pairwise momentum estimator. The data products (galaxy
  catalogs, CMB maps, estimator weights) are public through DESI and ACT
  data releases.

\item \textbf{Split into radial bins.} Re-analyze the pairwise velocity
  measurement in at least 3 radial bins:
  \begin{itemize}[nosep]
  \item Bin~A: $20 < r < 60\,h^{-1}$~Mpc (Newtonian regime, $\Geff \approx G_N$)
  \item Bin~B: $60 < r < 120\,h^{-1}$~Mpc (transition regime)
  \item Bin~C: $120 < r < 200\,h^{-1}$~Mpc (DFD-enhanced regime)
  \end{itemize}

\item \textbf{Compute $f\sigma_8$ in each bin.} Extract the effective
  growth rate from each radial bin independently. Under $\LCDM$, all
  bins yield the same $f\sigma_8$. Under DFD, Bin~C should be
  $5$--$15\%$ higher than Bin~A.

\item \textbf{DFD template.} For the analytic estimate:
  \begin{equation}
  \frac{v_{12}^{\rm DFD}(r)}{v_{12}^{\LCDM}(r)} = 1 + A\left(
  \frac{r}{r_\mu}\right)^2, \qquad r_\mu = \frac{2\pi}{k_\mu}
  \approx 600\,h^{-1}\,\text{Mpc}, \quad A \approx 0.04.
  \end{equation}

\item \textbf{Model comparison.} Fit:
  \begin{itemize}[nosep]
  \item Model~A ($\LCDM$): flat $f\sigma_8$ across all bins.
  \item Model~B (DFD): rising $f\sigma_8$ with the predicted
    $(r/r_\mu)^2$ shape.
  \end{itemize}
  Report $\Delta\chi^2$, posterior on the enhancement amplitude,
  and Bayesian evidence ratio.

\item \textbf{Systematic checks.} (a)~Verify that the scale split is
  not degenerate with galaxy bias evolution. (b)~Check for photo-$z$
  contamination at large separations. (c)~Compare with the ACT~DR6
  lensing-independent estimate.
\end{enumerate}

\section{Projected Sensitivity}

Current DESI~DR1 + ACT~DR6 precision: $\sim 10\%$ per radial bin at
$9.3\sigma$ total significance. DFD's predicted enhancement at
$r > 100\,h^{-1}$~Mpc is $\sim 10\%$, placing the signal at
$\sim 1\sigma$ with current data.

DESI~DR2 + Simons Observatory (2026--2027): $\sim 3\times$ improvement,
making this a $\sim 3\sigma$ test.


%=============================================================================
%=============================================================================
\part{Protocols 3--6: Unchanged from W4i14 (Summary Only)}
\label{part:unchanged}
%=============================================================================
%=============================================================================

The following protocols are unaffected by the distance-bias paradigm
correction because they test the scalar field coupling $|\lambda - 1|$
directly, not cosmological distance measures:

\begin{center}
\renewcommand{\arraystretch}{1.3}
\begin{tabular}{@{}clccl@{}}
\toprule
\textbf{Proto.} & \textbf{Experiment} & \textbf{Cost} &
  \textbf{Time} & \textbf{W4i15 Status} \\
\midrule
3 & GPS 59~ps archival (W4i14 Protocol~2) & \$50--100K & 6~mo &
  UNCHANGED \\
4 & LIGO O4 scalar mode (W4i14 Protocol~3) & \$100--500K & 12~mo &
  UNCHANGED \\
5 & SQMS Phase~I (W4i14 Protocol~4) & \$3.2M & 24~mo &
  UNCHANGED \\
6 & MEMS in-gap (W4i14 Protocol~5) & \$200--500K & 18~mo &
  UNCHANGED \\
\bottomrule
\end{tabular}
\end{center}

The full specifications (satellites, receivers, software, derivation
chains, systematic checks, decision gates) are in W4i14 Protocols~2--5.
They carry forward without modification.

\textbf{One update for SQMS (Protocol~5)}: the $\omega^2$ loss rate
scaling flagged in W4i14 Finding~5 remains \textbf{UNVERIFIED} from
the action principle. This is a blocker for the Q-ratio diagnostic
(0.49 vs 3.41) and should be resolved by Agent~1 before Phase~I
proposal submission.


%=============================================================================
%=============================================================================
\part{The \texttt{mochi\_class} Boltzmann Code: Promoted to Critical Path}
\label{part:boltzmann}
%=============================================================================
%=============================================================================

\begin{tcolorbox}[colback=red!5!white, colframe=warncolor,
  title=\textbf{CRITICAL PATH: \texttt{mochi\_class} Now Gates All
  Cosmological Predictions}]
\textbf{Budget}: \$40--80K\\
\textbf{Timeline}: 2--3 months\\
\textbf{Why critical}: With BAO $\approx$ $\LCDM$, all DFD-distinctive
cosmological signals come from $\Geff(k,a)$. The analytic estimates
(probe-dependent $S_8$, kSZ enhancement, $A_L$ scale dependence)
are order-of-magnitude only. Boltzmann-precision predictions are needed
to convert $\sim 1\sigma$ hints into $> 3\sigma$ tests.\\
\textbf{Deliverables}: $P(k,z)$ with DFD $\Geff$; $C_l^{TT/TE/EE}$;
$C_l^{\kappa\kappa}$; $S_8$ per probe window function; $f\sigma_8(k,z)$;
kSZ pairwise velocity template.
\end{tcolorbox}

\section{Why It Gates Everything}

\begin{enumerate}[nosep]
\item \textbf{$S_8$ probe hierarchy} (Protocol~1R): The analytic estimate
  gives the correct ordering but uncertain amplitudes.
  \texttt{mochi\_class} provides the exact $\Delta S_8$ per probe.

\item \textbf{kSZ pairwise velocities} (Protocol~2N): The
  $v_{12}^{\rm DFD}(r)/v_{12}^{\LCDM}(r) = 1 + A(r/r_\mu)^2$ template
  is approximate. The Boltzmann code computes the exact nonlinear
  $P(k)$ modification and its projection onto pairwise velocities.

\item \textbf{CMB lensing $A_L$}: DFD predicts $A_L \sim 1.10$--$1.20$
  with scale dependence ($A_L(l)$ varying across multipoles).
  Boltzmann-precision $C_l^{\kappa\kappa}$ is needed to quantify this.

\item \textbf{Neutrino mass}: DFD predicts $\Sigma m_\nu = 61.4$~meV.
  The margin with DESI~DR2 + Planck~PR4 is $\sim 2$~meV (i14). The
  Boltzmann code determines whether DFD shifts the $\Sigma m_\nu$
  posterior through $\Geff$-modified growth.

\item \textbf{$f\sigma_8$ scale-resolved growth}: DESI RSD measurements
  at $\sim 2$--$5\%$ per bin are marginally constraining for the
  DFD $3$--$5\%$ suppression at $z < 0.5$. Boltzmann precision is
  needed to determine the exact amplitude and $z$-dependence.
\end{enumerate}

\section{Implementation (Unchanged from W4i14 Protocol~6)}

\begin{enumerate}[label=\textbf{B.\arabic*:}, nosep]
\item Install \texttt{mochi\_class} from
  \texttt{github.com/music-of-cosmos/mochi\_class\_public}.
\item Implement DFD Horndeski mapping: $\alpha_M(a) = (d\psi_0/d\ln a)/(aH)$,
  $\alpha_B = \alpha_T = 0$.
\item Run CMB $C_l$, $P(k)$, $C_l^{\kappa\kappa}$, BAO distances.
\item Generate per-probe $S_8$ predictions.
\item Generate kSZ pairwise velocity template.
\item Compare with Planck~PR4 + DESI~DR2 + DES~Y6 + eROSITA.
\end{enumerate}


%=============================================================================
%=============================================================================
\part{Redesigned Programme: Budget, Timeline, Optimal Sequence}
\label{part:consolidated}
%=============================================================================
%=============================================================================

\section{Complete Experimental Programme (W4i15)}

\begin{center}
\renewcommand{\arraystretch}{1.3}
\begin{longtable}{@{}clcccl@{}}
\toprule
\textbf{Rank} & \textbf{Experiment} & \textbf{Cost} & \textbf{Time} &
  \textbf{Status} & \textbf{W4i15 Update} \\
\midrule
\endhead
0a & $S_8$ probe hierarchy & \$0 & 2--4 wks & \textbf{EXECUTE NOW} &
  \textbf{NEW}: replaces $D_H/r_d$ \\
0b & kSZ velocity scale split & \$0 & 2--4 wks & \textbf{EXECUTE NOW} &
  \textbf{NEW}: DESI+ACT data public \\
0c & Rotation curves & \$0 & 3--6 mo & Unchanged & --- \\
\midrule
$\infty$ & \texttt{mochi\_class} Boltzmann & \$40--80K & 2--3 mo &
  \textbf{CRITICAL PATH} & Gates all cosmo predictions \\
\midrule
1a & GPS 59~ps archival & \$50--100K & 6 mo & UNCHANGED & --- \\
1b & LIGO O4 scalar mode & \$100--500K & 12 mo & UNCHANGED & --- \\
\midrule
I & SQMS Phase~I & \$3.2M & 24 mo & UNCHANGED &
  $\omega^2$ loss rate still unverified \\
II & MEMS in-gap test & \$200--500K & 18 mo & UNCHANGED &
  $Q_{\rm mech}$ gap 150$\times$ \\
\midrule
\multicolumn{6}{l}{\textbf{KILLED (confirmed + new):}} \\
--- & $D_H/r_d$ sign change & --- & --- & \textbf{KILLED (i15)} &
  Distance-bias correction \\
--- & Torsion direct force & --- & --- & KILLED (i13) & $G/c^4$ \\
--- & Optical clock altitude & --- & --- & KILLED (i13) & $10^{-41}$ \\
--- & Neutron interferometry & --- & --- & KILLED (i13) & $10^{-22}$ rad \\
--- & Casimir scalar probe & --- & --- & KILLED (i13) & GR-standard \\
--- & Muon $g{-}2$ & --- & --- & NULL (i13) & $10^{-15}$ \\
--- & Dark SRF reinterpretation & --- & --- & INVALIDATED (i12) & \\
\bottomrule
\end{longtable}
\end{center}


\section{Optimal Spending Sequence}

\begin{tcolorbox}[colback=yellow!5!white, colframe=yellow!60!black,
  title=\textbf{Optimal Decision Cascade (W4i15)}]

\textbf{Phase 0a [\$0, Weeks 1--4]:}
Execute $S_8$ probe hierarchy and kSZ velocity scale split simultaneously.
Both use published data, standard software, and require only
analysis effort.\\[4pt]

\textbf{Decision Gate 0a:} If either test yields $> 2\sigma$ signal
consistent with DFD $\Geff(k)$: fast-track \texttt{mochi\_class}.
If both tests yield null results at $> 3\sigma$: re-evaluate whether
$k_\mu \sim 0.01\,h\,\text{Mpc}^{-1}$ is correct (possible recalibration
of crossover scale from RAR).\\[4pt]

\textbf{Phase 0b [\$40--80K, Months 1--3]:}
\texttt{mochi\_class} Boltzmann code. Generates Boltzmann-precision
predictions for all $\Geff$-dependent observables. Required regardless
of Phase~0a outcome --- even a null result needs quantification.\\[4pt]

\textbf{Phase 0c [\$50--100K, Months 1--6]:}
GPS 59~ps archival search. Independent of cosmological tests (probes
$|\lambda - 1|$ directly). Can run in parallel with Phase~0b.\\[4pt]

\textbf{Decision Gate 0c:} If GPS multi-GNSS ratio $R = 1.070 \pm 0.01$
detected at SNR~$> 5$: proceed to multi-GNSS campaign and SQMS.
If null: tighten bound on $|\lambda - 1|$ from GPS data.\\[4pt]

\textbf{Phase 1 [\$100--500K, Months 3--15]:}
LIGO O4 scalar breathing mode search. Start with O4a data (on GWOSC);
expand to O4b (May 2026) and O4c (December 2026) as released.\\[4pt]

\textbf{Phase 2 [\$3.2M, Months 6--30]:}
SQMS Phase~I. Proceed if $\omega^2$ loss rate derivation is verified by
Agent~1. The Q-ratio diagnostic (0.49 vs 3.41) is independent of the
cosmological framework.\\[4pt]

\textbf{Phase 3 [\$200--500K, Months 12--30]:}
MEMS in-gap test. Proceed when $Q_{\rm mech} = 10^9$ is demonstrated
cryogenically (currently 150$\times$ gap).\\[4pt]

\textbf{Phase 4 [2027--2028+]:}
Euclid~DR1, DESI~DR3, Simons Observatory, JUNO, CMB-S4.

\end{tcolorbox}


\section{Budget Summary}

\begin{center}
\renewcommand{\arraystretch}{1.2}
\begin{tabular}{@{}lrl@{}}
\toprule
\textbf{Phase} & \textbf{Cost} & \textbf{Tests} \\
\midrule
Phase 0a (zero-cost tests) & \$0 & $S_8$, kSZ, rotation curves \\
Phase 0b (Boltzmann code) & \$40--80K & All cosmo predictions \\
Phase 0c (GPS archival) & \$50--100K & $|\lambda - 1|$ direct \\
Phase 1 (LIGO O4) & \$100--500K & Scalar breathing mode \\
\midrule
\textbf{Total Phase 0+1} & \textbf{\$190--680K} & \\
\midrule
Phase 2 (SQMS) & \$3.2M & Q-ratio, absolute bound \\
Phase 3 (MEMS) & \$200--500K & $\Qpsi$ threshold \\
\midrule
\textbf{Total all phases} & \textbf{\$3.6--4.4M} & \\
\bottomrule
\end{tabular}
\end{center}


%=============================================================================
\section*{Derivation Chains (Updated)}
%=============================================================================

\begin{enumerate}
\item \textbf{$S_8$ probe hierarchy (NEW, \$0)}: Action (section02
  Eq.~2.6) $\to$ field equation with $\mu(x)=x/(1+x)$ $\to$ linearized
  perturbations $\to$ $\Geff(k) = G_N(k^2 + k_\mu^2)/k^2$ $\to$
  scale-dependent $P(k)$ $\to$ probe-dependent $S_8$ weighting $\to$
  ordered hierarchy eROSITA $>$ Planck $>$ KiDS $>$ DES.
  \textbf{Zero free parameters. Testable NOW.}

\item \textbf{kSZ pairwise velocity (NEW, \$0)}: Same $\Geff(k)$ $\to$
  scale-dependent $f(k)$ $\to$ $v_{12}(r)$ enhanced at large $r$ $\to$
  $+10\%$ at $r > 100\,h^{-1}$~Mpc.
  \textbf{Zero free parameters. Testable NOW with DESI+ACT.}

\item \textbf{GPS 59~ps}: Action $\to$ $\psi_{\rm grav}$ gradient $\to$
  nonreciprocal $\delta t = \Delta\psi/(2c) = 59$~ps $\to$
  multi-GNSS ratio $R = 1.070$. \textbf{Unchanged.}

\item \textbf{LIGO 6th channel}: Action $\to$ $\Box\psi$ $\to$
  scalar breathing $h_b$ $\to$ common-mode rejected $\to$
  reach $1.5\ee{-14}$ (O4). \textbf{Unchanged.}

\item \textbf{SQMS Q-ratio}: Action $\to$ EM-$\psi$ coupling $\to$
  $\Gamma_\psi \propto \omega^2$ $\to$ $R_Q = 0.49$.
  \textbf{Unchanged, but $\omega^2$ scaling UNVERIFIED.}

\item \textbf{MEMS in-gap}: Action $\to$ $\psi$-phonon coupling $\to$
  $\Qpsi = 2.2\ee{8}$ $\to$ saturation plateau.
  \textbf{Unchanged. Technical risk: $Q_{\rm mech}$ gap 150$\times$.}
\end{enumerate}


%=============================================================================
\section*{Summary of Changes from W4i14}
%=============================================================================

\begin{center}
\renewcommand{\arraystretch}{1.2}
\begin{tabular}{@{}lll@{}}
\toprule
\textbf{Item} & \textbf{W4i14} & \textbf{W4i15} \\
\midrule
$D_H/r_d$ sign-change & Protocol 1, EXECUTE NOW & \textbf{KILLED} \\
$S_8$ probe hierarchy & Mentioned in Cosmo & \textbf{Protocol 1R, \$0} \\
kSZ velocity split & Prediction 17 (New Pred.) & \textbf{Protocol 2N, \$0} \\
\texttt{mochi\_class} & Protocol 6, important & \textbf{CRITICAL PATH} \\
GPS 59~ps & Protocol 2 & Unchanged \\
LIGO O4 & Protocol 3 & Unchanged \\
SQMS & Protocol 4 & Unchanged \\
MEMS & Protocol 5 & Unchanged \\
Spending sequence & DESI $\to$ GPS $\to$ LIGO & $S_8$+kSZ $\to$ Boltzmann
  $\to$ GPS $\to$ LIGO \\
\bottomrule
\end{tabular}
\end{center}


\end{document}
